% ── sections/06_declarations.tex ─────────────────────────────────────────────

\section*{Declarations}

\subsection*{Study Limitations}

\textbf{Single-seed evaluation.} Most reported results are based on single training runs.
Multi-seed validation of the Optuna best configuration (hidden=32, $L$=4, seeds 1--5)
yields mean AUPR~=~0.742~$\pm$~0.008, comparable to the M1 baseline mean (0.743), indicating
that the single-trial Optuna result (0.8023) is seed-dependent. The BatchNorm ablation
result is controlled and reproducible. For the six-network model (AUPR~=~0.8067), full
multi-seed variance estimation remains as future work.

\textbf{Gain decomposition.} A control experiment (Benchmark GCN, all six networks)
achieves AUPR~=~0.7987, isolating the data contribution (+0.054) from the architecture
contribution (+0.008). The majority of the +5.9\% total gain derives from multi-network
data integration rather than the architectural improvements introduced in EMGNNImproved.

\textbf{Network weights circular dependency.} The learned per-network importance weights
(Section~4.3) assign the highest weight to CPDB (0.210). However, because the CPDB test
set is used as the optimisation target throughout training and evaluation, the model
receives a training signal that directly rewards features derived from CPDB. This creates
a circular dependency: the CPDB network may be up-weighted partly because it is the
evaluation network, not solely because of its biological quality. Independent evaluation
on a held-out network or cancer type would be required to confirm that the weight
rankings reflect intrinsic network quality.

\textbf{Zero-baseline IG attribution.} Integrated Gradients scores are computed relative
to a zero-feature baseline (see Section~\ref{sec:theory}). Features with high raw
magnitudes (such as un-centred methylation probes) may receive inflated attributions
independent of their true discriminative contribution. Baseline sensitivity analysis
using alternative reference points was not performed in this study.

\textbf{Validation split protocol.} The validation split is constructed from a 10\%
subset of training nodes rather than the official \texttt{mask\_val} provided in the HDF5
files, which may introduce minor protocol inconsistencies relative to the original
EMGNN evaluation.

\textbf{Normalisation ablation.} GraphNorm and LayerNorm are implemented as alternative
normalisation strategies but were not evaluated experimentally against the no-normalisation
baseline. The ablation study only contrasts BatchNorm1d with no normalisation; a full
three-way comparison is deferred to future work.

\textbf{Methodology~5 evaluation pending.} The eleven advanced techniques (Focal Loss,
DropEdge, heterophily-aware gating, GraphMAE, GPS, cross-network attention, hypergraph,
HIPGNN, positional encoding, PINNACLE, GNNExplainer) are fully implemented and
syntactically verified but have not been evaluated in a systematic ablation study on the
six-network configuration. Individual and combined performance gains remain to be
quantified. Three modules require external datasets (MSigDB GMT files, Laplacian
eigenvectors, PINNACLE embeddings from HuggingFace) that were not available during
this study.

\subsection*{Acknowledgements}

The author thanks the original EMGNN authors (Chatzianastasis et al.) and the EMOGI
benchmark authors (Schulte-Sasse et al.) for releasing their code and datasets publicly,
which made this reproducibility and extension study possible.

\subsection*{Funding Source}

None. This work was conducted as part of a final-year undergraduate capstone project at
the National University of Singapore and received no external funding.

\subsection*{Competing Interests}

The author declares no competing interests in relation to this publication.

\subsection*{Human and Animal Related Study}

This study involves computational analysis of publicly available genomic datasets (TCGA)
and PPI network databases. No human subjects were recruited and no animals were used.
Ethical approval and informed consent are therefore not applicable.

